\documentclass[a4paper,12pt]{article}
\usepackage{amsmath,amssymb}
\usepackage{xcolor}
\usepackage{tcolorbox}
\usepackage{geometry}
\geometry{left=2cm,right=2cm,top=2cm,bottom=2cm}

% --- boxes ---------------------------------------------------
\definecolor{SolBG}{rgb}{0.95,0.98,1.00}
\definecolor{ExeBG}{rgb}{0.93,0.93,1.00}

\newtcolorbox{exercice}[1][]{
  colback=ExeBG,
  colframe=purple!70,
  title=\textbf{#1},
  fonttitle=\bfseries,
}

\newtcolorbox{solution}[1][]{
  colback=SolBG,
  colframe=blue!70!black,
  title=\textbf{#1},
  fonttitle=\bfseries,
}

%--------------------------------------------------------------
\begin{document}

%%%%%%%%%%%%%%%%%%%%  Exercice 13  %%%%%%%%%%%%%%%%%%%%%%%%%%%%%
\begin{exercice}[Exercice 13 – Calculer]
Dans chaque cas on détermine (à 0,1 près) les longueurs des trois côtés
et les mesures des trois angles du triangle $ABC$.
\end{exercice}

\begin{solution}[Correction]
\textbf{Rappels.}  
Loi des cosinus (Al-Kashi)
\(
BC^{2}=AB^{2}+AC^{2}-2\,AB\cdot AC\cos\widehat A
\)
et loi des sinus
\(
\displaystyle\frac{AB}{\sin\widehat C}
      =\frac{BC}{\sin\widehat A}
      =\frac{AC}{\sin\widehat B}.
\)

\begin{enumerate}
\item[\bfseries 1.] ${AB}=3,\;AC=7,\;\widehat A=37^{\circ}$  
  \[
  BC^{2}=3^{2}+7^{2}-2\!\times\!3\!\times\!7\cos37^{\circ}\approx24.5
  \Longrightarrow BC\approx4.9.
  \]
  Dès lors
  \(
  \sin\widehat B
  =\dfrac{AC\,\sin\widehat A}{BC}
  \approx\dfrac{7\times0.6018}{4.9}\approx0.853
  \Rightarrow
  \widehat B\approx58.6^{\circ},
  \)
  d’où
  \(
  \widehat C=180^{\circ}-37^{\circ}-58.6^{\circ}\approx84.4^{\circ}.
  \)

\smallskip
\item[\bfseries 2.] ${AB}=1,\;AC=4,\;\widehat B=110^{\circ}$  
  Utilisons Al-Kashi à l’angle $B$ :
  \[
  AC^{2}=AB^{2}+BC^{2}-2\,AB\cdot BC\cos\widehat B
  \Longrightarrow
  BC^{2}+0.684\,BC-15=0
  \]
  (car $\cos110^{\circ}\simeq-0.342$).  
  La solution positive est $BC\approx3.5$.  
  Loi des sinus :
  \(
  \sin\widehat C
  =\dfrac{AB\,\sin110^{\circ}}{AC}\approx0.235
  \Rightarrow\widehat C\approx13.6^{\circ},
  \)
  puis $\widehat A\approx56.4^{\circ}$.

\smallskip
\item[\bfseries 3.] ${BC}=9,\;\widehat B=72^{\circ},\;\widehat C=34^{\circ}$  
  $\widehat A=180^{\circ}-72^{\circ}-34^{\circ}=74^{\circ}$.  
  Par la loi des sinus
  \[
  AB=\frac{\sin34^{\circ}}{\sin74^{\circ}}\;9\approx5.2,
  \qquad
  AC=\frac{\sin72^{\circ}}{\sin74^{\circ}}\;9\approx8.9.
  \]

\smallskip
\item[\bfseries 4.] ${BC}=11,\;AC=3,\;\widehat B=72^{\circ}$  
  Ici
  \(
  \sin\widehat A=\dfrac{BC\sin\widehat B}{AC}
  =\dfrac{11\times0.951}{3}=3.49>1.
  \)
  L’équation $\sin\widehat A>1$ est impossible : \textbf{aucun triangle
  ne répond aux données}.  

\smallskip
\item[\bfseries 5.] ${AC}=17,\;\widehat B=64^{\circ},\;\widehat C=71^{\circ}$  
  $\widehat A=45^{\circ}$ et $k=AC/\sin\widehat B
  =17/\sin64^{\circ}\approx18.9$.  
  Alors
  \[
    BC=k\sin45^{\circ}\approx13.4,
    \qquad
    AB=k\sin71^{\circ}\approx17.9.
  \]
\end{enumerate}
\end{solution}

%%%%%%%%%%%%%%%%%%%%%%%%%%%%%%%%%%%%%%%%%%%%%%%%%%%%%%%%%%%%%%%
% Exercice 18
%%%%%%%%%%%%%%%%%%%%%%%%%%%%%%%%%%%%%%%%%%%%%%%%%%%%%%%%%%%%%%%
\begin{exercice}[Exercice 18 – Chercher / Représenter / Calculer]
Deux points $A$ et $B$ sont tels que $AB=4$.
Déterminer l’ensemble des points $M$ pour lesquels
$\overrightarrow{AM}\cdot\overrightarrow{AB}=10$.
\end{exercice}

\begin{solution}[Correction]
Plaçons un repère où $A$ est l’origine et $B$ est sur l’axe $x$ :
$B(4,0)$, donc $\vec{AB}=(4,0)$.
Écrivons $M(x,y)$ avec $\vec{AM}=(x,y)$ :

\[
\vec{AM}\cdot\vec{AB}=4x=10\quad\Longrightarrow\quad x=2.5.
\]

L’ensemble recherché est la \textbf{droite} parallèle à $AB$,
d’équation $x=2.5$, c’est-à-dire la droite perpendiculaire à $AB$ passant
à 2.5 cm du point $A$.
\end{solution}

%%%%%%%%%%%%%%%%%%%%%%%%%%%%%%%%%%%%%%%%%%%%%%%%%%%%%%%%%%%%%%%
% Exercice 19
%%%%%%%%%%%%%%%%%%%%%%%%%%%%%%%%%%%%%%%%%%%%%%%%%%%%%%%%%%%%%%%
\begin{exercice}[Exercice 19 – Chercher / Représenter / Calculer]
Deux points $A$ et $B$ sont tels que $AB=1$.
Déterminer l’ensemble des points $M$ pour lesquels
$\overrightarrow{MA}\cdot\overrightarrow{MB}=\dfrac32$.
\end{exercice}

\begin{solution}[Correction]
Toujours $A$ à l’origine, $B(1,0)$, donc
$\vec{MB}=(x-1,y)$ et $\vec{MA}=(x,y)$.
\[
\vec{MA}\cdot\vec{MB}=x(x-1)+y^2=\frac32.
\]
Développons :
\(
x^{2}-x + y^{2}=1.5.
\)
C’est le cercle de centre $\bigl(\tfrac12,0\bigr)$ et de rayon
$r=\sqrt{1.5+\tfrac14}= \sqrt{1.75}\simeq1.32$.
\end{solution}

%%%%%%%%%%%%%%%%%%%%%%%%%%%%%%%%%%%%%%%%%%%%%%%%%%%%%%%%%%%%%%%
% Exercice 20
%%%%%%%%%%%%%%%%%%%%%%%%%%%%%%%%%%%%%%%%%%%%%%%%%%%%%%%%%%%%%%%
\begin{exercice}[Exercice 20 – Chercher / Représenter / Calculer]
Deux points $A$ et $B$ sont tels que $AB=5$ cm.
Déterminer l’ensemble des points $M$ pour lesquels l’aire du triangle
$ABM$ est $3$ cm$^{2}$.
\end{exercice}

\begin{solution}[Correction]
L’aire d’un triangle vaut
\(
\mathcal A=\dfrac12\,AB\cdot h,
\)
où $h$ est la distance de $M$ à la droite $(AB)$.
Ici
\(
3=\dfrac12\times5\times h\Longrightarrow h=\dfrac{6}{5}=1.2\text{ cm}.
\)
Le lieu cherché est donc constitué des \textbf{deux droites parallèles}
à $(AB)$ situées à 1.2 cm de $(AB)$ (une de chaque côté).
\end{solution}

%%%%%%%%%%%%%%%%%%%%%%%%%%%%%%%%%%%%%%%%%%%%%%%%%%%%%%%%%%%%%%%
% Exercice 21
%%%%%%%%%%%%%%%%%%%%%%%%%%%%%%%%%%%%%%%%%%%%%%%%%%%%%%%%%%%%%%%
\begin{exercice}[Exercice 21 – Représenter / Calculer]
Soient $A(1;3)$ et $B(-2;5)$.
Déterminer l’ensemble des points $M(x,y)$ tels que
les vecteurs $2\overrightarrow{MA}+\overrightarrow{MB}$ et
$\overrightarrow{MA}+2\overrightarrow{MB}$ soient orthogonaux.
\end{exercice}

\begin{solution}[Correction]

\textbf{1. Mise en équation}

\[
\vec{MA}=(1-x,\;3-y),\qquad
\vec{MB}=(-2-x,\;5-y).
\]

Posons
\[
\vec v = 2\vec{MA}+\vec{MB},\qquad
\vec w = \vec{MA}+2\vec{MB}.
\]
On veut $\vec v\cdot\vec w=0$.

\[
\vec v = (\,3-3x,\; 6-3y- y)=\bigl(3-3x,\;6-4y\bigr),
\]
\[
\vec w = (\,3-3x,\;3- y-2y) =\bigl(3-3x,\;3-3y\bigr).
\]

Produit scalaire :
\[
(3-3x)^2 +(6-4y)(3-3y)=0.
\]

Après développement et division par $9$ on obtient
\[
x^{2}+x+y^{2}-8y+\frac{143}{9}=0.
\]

\textbf{2. Centre et rayon}

Complétons les carrés :
\[
(x+\tfrac12)^{2}+\bigl(y-4\bigr)^{2}=\frac{29}{81}\approx0.36.
\]

\textbf{3. Lieu géométrique}

Il s’agit d’un \textbf{cercle} de centre
$\displaystyle C\!\Bigl(-\tfrac12,\,4\Bigr)$
et de rayon
\(
r=\sqrt{29}/9\approx0.6.
\)
\end{solution}

\end{document}
